% OCEANS 2027 full paper (IEEE conference format).
% Status: draft started 2026-11-09, the day of pool session S1, before any pool recording has been scored.
% Sections II (Methods) and III (Real-data benchmark and S1 protocol) are drafted; the others are stubs.
% Every number is taken from a repository source (RESULTS_E7_SIM_DRYRUN.md, RESEARCH_PLAN.md,
% docs/RECORDING_SESSION_CHECKLIST.md, sessions/2026-11-09_s1/, configs/, uwse/, scripts/,
% data/bench/v0_s1/). Real-data quantities that do not exist yet are marked \tbd.
% References missing from the abstract's bibliography are marked [CITE: ...] and must be resolved
% before submission; no new bibliography entries have been added.
% Build (from the uw_speech_enhancement root):
%   pdflatex -interaction=nonstopmode -output-directory paper/build paper/oceans2027_full.tex   (twice)
\documentclass[conference]{IEEEtran}
\usepackage[T1]{fontenc}
\usepackage{amsmath}
\usepackage{graphicx}
\usepackage{cite}
\usepackage{url}

\newcommand{\tbd}{\textbf{[TBD]}}
\newcommand{\citeph}[1]{[CITE: #1]}

\begin{document}

\title{Speech Enhancement versus Recognizer Adaptation for Underwater Speech Recognition over a Projector-to-Hydrophone Link}

% Sole author (lead's decision, 2026-11-19).
\author{\IEEEauthorblockN{Ryota Sato}
\IEEEauthorblockA{Department of Electrical Engineering, Stanford University\\
Stanford, CA, USA\\
\texttt{ryos17@stanford.edu}}}

\maketitle

% ---------------------------------------------------------------------------------------------
\section{Introduction}
% TODO: write the introduction (motivation as in the abstract; enhance vs adapt; contributions).
\textit{TODO: Introduction.}

\section*{Related Work}
% TODO: decide whether Related Work is a numbered section or folded into the Introduction.
\textit{TODO: Related work (through-water speech \cite{nissen2017,roehrdanz2025,roehrdanz2026}; enhancement artefacts and ASR \cite{iwamoto2022}).}

% ---------------------------------------------------------------------------------------------
\section{Methods}
\label{sec:methods}

\subsection{Signal chain and simulator}
\label{sec:sim}
A clean utterance $s$ reaches the recorder through a drive low-pass filter (LPF), the projector, the
pool, the hydrophone and the converter. The channel simulator \texttt{uwsim} models this path as
\begin{equation}
y = Q\Big(\mathrm{HPF}\big(h_{\mathrm{hyd}} * \big(\textstyle\sum_{k=1}^{3} h_k * (\mathrm{LPF}_{8\,\mathrm{kHz}}\, s)^k * h_{\mathrm{ch}} + n\big)\big)\Big),
\label{eq:chain}
\end{equation}
where $(\cdot)^k$ is a pointwise power. Noise $n$ is added as acoustic pressure in the water, so it
passes through the hydrophone response $h_{\mathrm{hyd}}$, the preamplifier high-pass filter (HPF) and
the converter $Q$ together with the speech. The components are as follows.

\textit{Drive and projector.} The drive is low-passed at 8~kHz by a causal eighth-order Butterworth
filter (the playback bench instead applies an eighth-order design forward and backward,
Section~\ref{sec:benchbuild}). The projector is a third-order Hammerstein model whose linear term $h_1$
is the projector finite impulse response (FIR), an identity filter in v0 that is to be estimated from
the measured chain response; a piston directivity of radius 0.117~m weights the departure angle.

\textit{Channel $h_{\mathrm{ch}}$.} An image-source model (ISM) \cite{allen1979} of a
shoebox pool with a tilted floor gives the early field up to a crossover time $t_e$, drawn per scene from 30--50~ms (nominal
40~ms). The surface reflection
coefficient is real and signed, $R_s \in [-1.0, -0.9]$, and walls and floor reflect as a fluid--elastic
half-space evaluated at each image's incidence angle \cite{brekhovskikh1980}. A stochastic late tail (Gaussian noise with a per-octave exponential
decay) follows, energy-matched to the ISM decay at $t_e$. Below the crossover
$f_x = \max(c/4H, 500~\mathrm{Hz})$, with $c$ the sound speed and $H$ the maximum pool depth, a
modal sum replaces the image-source field. A time scaling by $(1+\varepsilon)$, $|\varepsilon| \le
10^{-3}$, models slow channel variation.

\textit{Noise.} In this work $n$ contains Minnaert bubble bursts \cite{minnaert1933}
($f_0 \approx 3.29/a$~Hz with the radius $a$ in metres, about 1.1--6.6~kHz for $a = 0.5$--$3$~mm;
rate up to 20 per second) and, in the
\emph{stat} training configuration below, stationary coloured Gaussian noise. No recorded noise bank
exists yet, so thruster and ambient noise are absent from the simulator. The noise is scaled to a
target speech-band signal-to-noise ratio (SNR), measured over 500--8000~Hz on speech-active samples,
which an energy voice activity detector (VAD) selects as the 20~ms frames within 40~dB of the loudest.

\textit{Receiver.} The receiver applies a hydrophone response, a 100~Hz second-order preamplifier
HPF and a 24-bit analog-to-digital converter (ADC) with a noise floor of $-110$~dB relative to full
scale (dBFS) and a hard clip at 0~dBFS. The
simulator runs at 48~kHz and resamples its output to 16~kHz (48~kHz for DeepFilterNet3).

\textit{Targets.} The linear chain response is $h_{\mathrm{chain}} = \mathrm{LPF}_{8\,\mathrm{kHz}}
* h_1 * h_{\mathrm{hyd}} * \mathrm{HPF}$, with no pool path, directivity or nonlinearity. The
enhancement target is the chain reference $R_{1,\mathrm{sim}} = s * h_{\mathrm{chain}}$, which carries
the projector and hydrophone coloration but no pool path; it is defined identically to the real-data
reference $R_1$ (Section~\ref{sec:refs}). The end-to-end target is $R_0 = s$. Both targets are
band-limited to $[f_{\mathrm{lo}}, 8~\mathrm{kHz}]$, aligned to the direct-path arrival and given the
same broadband gain fit against the received speech.

\textit{Calibration status.} Every parameter of \texttt{uwsim} v0 is an uncalibrated placeholder: the
pool takes the nominal dimensions of a competition pool (37 by 20~m, 3.35--4.27~m deep, not survey
data), the per-octave reverberation times are priors (1.0~s at 125~Hz falling to 0.4~s at 8~kHz, scaled
by a factor drawn from 0.7--1.3), and the transducer responses are identity filters. The full paper
replaces them with values fitted to pool measurements \tbd.

\subsection{Simulated test set}
\label{sec:simtest}
The simulated test set renders the v0-test selection, 200 LibriSpeech test-clean
utterances~\cite{panayotov2015} (5 from each of 40 speakers), with projector and hydrophone at 1~m
depth, on-axis, at ranges of 1, 3 and 6~m, using one fixed impulse response (IR) per range. Three
conditions are rendered: \emph{quiet} (30~dB SNR); \emph{noise0} (0~dB SNR, stationary Gaussian noise
with a $-3$~dB/octave slope); and \emph{bubbles0} (0~dB SNR, Minnaert bubble trains at 10 per second,
radii 0.5--3~mm). This gives nine condition-by-range cells of 200 utterances each.

\subsection{Systems}
\label{sec:systems}
All systems share one frozen recognizer, Whisper large-v3~\cite{radford2023}, except C8, which adapts
it. Labels follow the project's system list.
\begin{itemize}
\item \textbf{C0}: unprocessed hydrophone audio.
\item \textbf{C1}: published speech enhancement (SE) front ends used zero-shot with their released
weights: single-channel weighted prediction error (WPE) dereverberation \cite{nakatani2010,drude2018},
MetricGAN+ \cite{fu2021}, SepFormer \cite{subakan2021} trained on WHAM! (16~kHz) and on WHAMR!
(8~kHz) \cite{wichern2019,maciejewski2020}, and DeepFilterNet3 (DFN3)~\cite{schroeter2023}.
\item \textbf{C2}: DFN3 fully fine-tuned on \texttt{uwsim} toward $R_{1,\mathrm{sim}}$ (one seed,
50k steps, about $3.2\times10^{6}$ audio-seconds), fed the 48~kHz simulator output.
\item \textbf{C10}: a U-Net following the architecture and target of~\cite{roehrdanz2026} (four
levels, ideal-ratio-mask target, 0--6~kHz log-magnitude input) with our own optimizer settings, not the
published ones, trained from scratch on \texttt{uwsim} (one seed; early-stopped at step 24k).
\item \textbf{C8}: recognizer adaptation, Whisper large-v3 with low-rank adaptation
(LoRA)~\cite{hu2022}, decoding the unprocessed audio with no front end (Section~\ref{sec:c8}).
\item \textbf{Oracles}: Whisper on $R_1$ (chain coloration only) and on the clean digital source.
\end{itemize}
C2 and C10 were trained with no recorded noise bank, so their training noise is almost entirely
synthetic bubble trains. Front-end outputs are decoded by the same recognizer and decoding policy as C0.

\subsection{Recognizer adaptation and checkpoint selection}
\label{sec:c8}
\textit{Adapter.} LoRA adapters of rank 8 ($\alpha = 32$, dropout 0.05) are inserted on the attention
query and value projections of Whisper large-v3, following the LoRA setting of \cite{rong2025} (who adapt Whisper tiny to small). They are trained with cross-entropy (CE) on the LibriSpeech transcripts of simulated audio:
train-clean-100 utterances of at most 30~s, rendered through \texttt{uwsim} on the fly. Training uses
AdamW with $\beta = (0.9, 0.999)$, no weight decay, a peak learning rate of $10^{-4}$ after 1000
warm-up steps, gradient clipping at 5.0, about 64~s of audio per optimizer step and at most 30k steps.
Validation CE on 200 dev-clean utterances is computed every 2000 steps (patience 5), and a checkpoint
is kept every 4000 steps.

\textit{Training noise.} Each noise component is normalized to unit RMS and weighted by an
independent level, bubbles (rate drawn from 0--20 per second) by $-20$ to 0~dB, and the sum is then
scaled to the target SNR. In the \emph{stat} configuration, half of the training scenes also add
stationary coloured Gaussian noise weighted by $-10$ to 0~dB, so its level relative to the bubbles
spans $-10$ to $+20$~dB, and it is present even in scenes without bubbles. Its spectral slopes are drawn
from $[-6,-4] \cup [-2,0]$~dB/octave. The band $(-4,-2)$~dB/octave around the noise0 test slope is held
out, so the test noise is unseen in spectral tilt, although it comes from the same generator family.
Training SNRs are drawn from $-10$ to 25~dB. Adapters trained on bubble noise alone, and a second
\emph{stat} seed, serve as simulation references only.

\textit{Checkpoint selection.} Validation CE did not rank adapters by word error rate (WER) in our
runs, so the checkpoint is chosen by WER on a simulated development set that shares no utterance,
speaker or scene geometry with the test set: 120 LibriSpeech dev-clean utterances at ranges of 1.5, 3.5
and 5.5~m, 1.5~m depth, with the receiver at 0.35 of the pool width (generator seed 20261104). Four
conditions are rendered: quiet, bubbles at 0~dB SNR, and stationary noise at 0~dB SNR with slopes of
$-5$ and $-1$~dB/octave. Both slopes lie inside the training bands, so the held-out test band never
influences selection. The criterion is the lowest mean of the four per-condition corpus WERs under the
guarded decoding of Section~\ref{sec:decode}; every kept checkpoint and the final one are candidates.

\textit{Frozen adapter.} The adapter carried forward to real data is \emph{stat} seed 1 at step 30k
(SHA-256 \texttt{5063f02f\ldots f52f90}), the adapter fixed by the protocol of 2026-11-05 as the
lowest development-WER checkpoint of \emph{stat} seed 1. Its development-set mean WER is 17.64\%,
against 22.65\% for the unadapted recognizer. An exploratory average of the two \emph{stat} seeds,
selected by the same rule (17.69\%), was not carried forward. For scoring, the adapter is merged into
the base weights in single precision (fp32) before the cast to half precision (fp16).

\subsection{Guarded decoding}
\label{sec:decode}
Every system is decoded with \texttt{openai-whisper} (version 20250625, weights SHA-256 \texttt{e5b1a55b\ldots}) in fp16,
in English, with the transcription task, no prompt and no previous-text conditioning. Before decoding,
each cut is normalized to a root-mean-square (RMS) level of $-26$~dBFS over its speech-active region, which a VAD on the aligned
$R_0$ determines; the same cut feeds every system.

Adapted recognizers can transcribe an utterance correctly and then fail to emit the end-of-text token,
appending long invented text. The decoding policy therefore guards every system identically:
\begin{itemize}
\item beam search of width 5 at temperature 0;
\item Whisper's temperature fallback (0.2, 0.4, 0.6, 0.8, 1.0 with best-of-5 sampling), triggered when
the compression ratio exceeds 2.4 or the mean log-probability falls below $-1.0$, with the sampling
seed fixed at 0 on every call;
\item a decode-length cap of $\min(\lceil 8~\mathrm{tokens/s} \times d\rceil + 16,\ 224)$ tokens for a
cut of duration $d$. Correct transcripts of the simulated clips run at a median of 3.7 and a maximum of
6.2 tokens per second, below the 8 tokens per second of the cap.
\end{itemize}
The unguarded policy (temperature 0, no cap) is reported only as a sensitivity check.

\subsection{Evaluation metrics and statistics}
\label{sec:stats}
The evaluation protocol was frozen on 2026-11-05, with additions on 2026-11-08 (the substitution
flag, two-band SNR and sex reporting, and the per-utterance SNR method) and the S1 amendment of
2026-11-09 (Section~\ref{sec:endpoint}), all made before any pool recording existed and so before any
was scored. It was fixed after the simulation results were seen but before any real recording existed,
so it is pre-registered with respect to the real data; any later change is reported as a deviation.

\textit{WER.} Whisper's English text normalizer~\cite{radford2023} is applied to hypothesis and
reference, and corpus WER is $(S+D+I)/N_{\mathrm{ref}}$, the summed substitutions, deletions and
insertions over the summed reference words.

\textit{Confidence intervals.} Corpus WER carries a 95\% confidence interval (CI) from a percentile
bootstrap over speaker clusters, resampled with replacement within position-by-thruster strata
(10\,000 replicates); a speaker in two strata forms two clusters. System differences are paired:
$\Delta\mathrm{WER}$ between two systems uses a paired speaker-cluster bootstrap with 2000 resamples,
drawing the same resampled clusters for both systems. A difference is called significant when its 95\%
CI excludes zero.
% NOTE (lead, 2026-11-10): scripts/score.py uses one --n-boot for all bootstraps, so S1 is scored in two
% runs: --n-boot 10000 for corpus CIs, --n-boot 2000 for paired dWER. No code change.
These CIs cover the recorded positions in the recorded session only.

\textit{Hallucination flags.} Two flags are reported per system. The first marks an utterance whose
hypothesis-to-reference length ratio exceeds 2 or whose compression ratio exceeds 2.4. The second, a
substitution flag, marks an utterance with $n_{\mathrm{hyp}} \ge 0.5\,n_{\mathrm{ref}}$ and
$\mathrm{hits}/n_{\mathrm{ref}} < 0.3$, both counted after normalization. ``WER without flagged'' is
computed on a single paired subset that drops an utterance flagged for any compared system. The flags
never filter the primary endpoint, which is WER on all utterances. Reference-free confidence scores are
not used as filters, because Whisper's mean log-probability is biased between C0 and C8. A runaway loop
is counted when an utterance's insertions reach $\max(10, 2\times$ its reference length$)$.

\textit{SNR and speaker sex (reporting only).} Real-data results report the speech-band SNR over both
50--8000~Hz and 500--8000~Hz, with the same band and energy-VAD definition as the simulator, applied
to the recorded thruster-off cut (speech plus ambient noise) rather than to the clean source. WER is also reported by speaker sex and
read against the median SNR of each sex in both bands; the female-minus-male gap carries a 95\% CI
from resampling speakers within sex.

\textit{Intrusive metrics.} Extended short-time objective intelligibility (ESTOI)~\cite{jensen2016} is
computed against $R_1$ at 16~kHz where $R_1$ exists (Section~\ref{sec:refs}).

% ---------------------------------------------------------------------------------------------
\section{Real-Data Benchmark and S1 Protocol}
\label{sec:bench}
This section describes pool session S1 (2026-11-09) as planned and pre-registered. No S1 recording
had been scored when it was written; every quantity measured at the session is marked \tbd.

\subsection{Bench construction and freezing}
\label{sec:benchbuild}
The playback bench is built from two utterance selections: v0-test (the 200 test-clean utterances of
Section~\ref{sec:simtest}) and v0-dev (40 dev-clean utterances). Each utterance is resampled to
48~kHz, RMS-normalized to $-26$~dBFS (peak capped at $-1$~dBFS, with any gain cut logged), given 5~ms
raised-cosine fades and passed through a zero-phase 8~kHz drive LPF (an eighth-order Butterworth
applied forward and backward, $-6$~dB at 8~kHz). Utterances are packed
into 28 content blocks of about 70~s (54--76~s; 23 from v0-test and 5 from v0-dev, 6 to 12 utterances
each),
each laid out as
\begin{center}
\small[0.5\,s silence][chirp][0.5\,s][utt$_1$][1.0\,s]$\cdots$[utt$_N$][0.5\,s][chirp][2.0\,s silence].
\end{center}
The sync chirp is a 0.25~s linear frequency-modulated (LFM) sweep from $f_{\mathrm{lo}} = 500$~Hz
(the value in the frozen S1 bench) to 8~kHz, with a Tukey (0.1) window and a
peak of $-20$~dBFS. It bypasses the drive LPF, so the matched-filter template is exactly what is
played. Silences and chirps take 16.3\% of block time.

For S1 the bench is built into \texttt{data/bench/v0\_s1} for position \texttt{p3} only, with the
default settings of the playlist and session-audio scripts. It was frozen before the session (built
2026-11-09 00:09): \texttt{SHA256SUMS.txt} lists all 68 files and itself hashes to
\texttt{6e6f6d02\ldots c5296357}; \texttt{schedule.json} hashes to \texttt{c150c540\ldots 6dc971} and
the speech-take timing file \texttt{session/p3\_speech.json} to \texttt{58d7e4cc\ldots 038009}. The
bench is never rebuilt after a take, because the import reads the same schedule and block timing files
that produced the played audio. The three-position bench \texttt{data/bench/v0} is kept for the
later session S2.

\subsection{Playback chain}
\label{sec:playback}
A laptop on battery plays each session file through a two-input, two-output USB interface whose first
output drives a battery-powered amplifier and the projector \tbd{} (projector and amplifier models). The
amplifier gain is set once per session, logged, and not touched again. Before the first take, a live
level check plays 15~s of the speech file and measures the drive on the loopback input; it passes when
the measured drive is within about 1~dB of the estimate, no input clips, and the hydrophone peaks keep
at least 6~dB of headroom. The nominal speech-active RMS source level is 150--160~dB re
1~$\mu$Pa at 1~m, an uncalibrated estimate from the projector's nominal sensitivity; the session values
of amplifier gain, interface output voltage and source level are \tbd.

\subsection{Recording rig and channel map}
\label{sec:rig}
Three recorders each write one file per take (Table~\ref{tab:chmap}). The USB interface records at
48~kHz the amplifier output through a 40~dB attenuator (loopback) and a monitor hydrophone at least
0.5~m from the projector, 90$^\circ$ off-axis. Two Daisy Seed boards each record two hydrophone
channels at 96~kHz and 24~bits. Each hydrophone keeps one channel identifier, \texttt{ch0} to
\texttt{ch3}, for the whole benchmark; \texttt{ch0} is the scored channel. The hydrophone elements,
preamplifier gains and recorder gains used at S1 are \tbd. Mounts are flooded or perforated, with no
air-filled tubing near the hydrophones.

\begin{table}[t]
\caption{S1 recorder channel map.}
\label{tab:chmap}
\centering
\small
\begin{tabular}{llll}
\hline
Recorder & Rate & File ch. & Role \\
\hline
USB interface & 48~kHz & 0 & loopback, 40~dB atten. \\
              &        & 1 & monitor hydrophone \\
Daisy A & 96~kHz & 0, 1 & \texttt{ch0} (scored), \texttt{ch1} \\
Daisy B & 96~kHz & 0, 1 & \texttt{ch2}, \texttt{ch3} \\
\hline
\end{tabular}
\end{table}

\subsection{Synchronization, clock-drift correction and import}
\label{sec:sync}
Every block carries its own start and end chirp, so each block is aligned on its own and the recorders
need no common clock. A recorder that records the loopback is aligned on it, which is exact. A recorder
without one, here both Daisy boards, is aligned on its first hydrophone's chirps, which carry the
acoustic delay. Recordings at 96~kHz are first resampled to the 48~kHz playback rate by an exact
rational ratio; raw sample indices in the output refer to the original file. The import proceeds in
four steps.
\begin{enumerate}
\item \textit{Coarse offset.} The matched-filter envelope of the sync channel, max-pooled to 10~ms, is
correlated with an impulse train at the nominal chirp onsets of the whole file, each impulse widened by
the largest drift expected over the take (200~ppm), so one missed chirp does not move the offset.
\item \textit{Per-block chirps.} Each chirp is searched within $\pm 0.5$~s of its prediction. A chirp
that is missing, too far from its prediction or below a peak-to-sidelobe ratio (PSR) of 6~dB is
replaced by the prediction of a take-level clock fit, and the block is flagged.
\item \textit{Clock drift.} A block's drift comes from its own two chirps when both are good, and
otherwise from a least-squares fit of chirp onset against nominal time over all good chirps of the
take.
\item \textit{Cuts.} Each channel is resampled onto the playback clock, anchored at the block's first
chirp, and every utterance is cut with 0.25~s of padding on either side.
\end{enumerate}
On synthetic takes with clock errors up to $\pm150$~ppm, the offset is recovered within 2 samples at
48~kHz and the cuts within 2--3 samples. In an end-to-end dry run on a synthetic session with recorder
clock errors of $+40$ and $-25$~ppm, every chirp passed quality control (QC) and both drifts were recovered exactly. A
Daisy frame drop removes 32 or 64 samples (0.3--0.7~ms); a drop between blocks is absorbed, and a drop
inside a block shifts its cuts by at most the dropped length, which is small against the padding. At scoring time, the
fine lag of each cut against the played utterance ($R_0$), within $\pm20$~ms, is estimated and applied
to the $R_0$ used for level normalization.

\subsection{S1 take plan}
\label{sec:takes}
S1 has 2~h of pool time and records one position, \texttt{p3}: a 3~m range with projector and
hydrophones at 1~m depth, on-axis. Table~\ref{tab:takes} lists the takes in order.

\textit{Noise-only takes.} Two 60~s noise-only takes open the position, one with the thruster off and
one with it on at the speech-take throttle (\tbd). For the thruster-on take, the thruster is at steady
speed before the recorders start and is stopped only after they stop.

\textit{Sweeps.} One sweep take plays a sync chirp and four exponential sine sweeps (ESS)
\cite{farina2000} from 150~Hz to 20~kHz, each 10.01~s long and followed by 5~s of
silence, at a peak of $-20$~dBFS. A low-frequency pre-shaping attenuates the sweep by about 12~dB well
below 300~Hz so that the projector is not overdriven; deconvolution removes it from the linear IR.

\textit{Speech (ABBA).} The 66.5~min speech take plays all 28 content blocks twice, in both thruster
states and adjacent in time: content block $j$ is played thruster-off then thruster-on for even $j$ and
in the reverse order for odd $j$, so the states read off, on, on, off, off, \ldots (ABBA). This gives
56 blocks, 112 chirps and 480 utterance presentations, and it removes the confound between thruster
state and time. The thruster (unit A, guarded frame) is switched at printed cues that fall right after a
block's end chirp, leaving about 3.5~s of spin-up before the next block.

\textit{Anchor.} After the speech take, the session's first block (9 utterances, thruster off) is
played again as a drift anchor.

\textit{Noise bank.} Four 90~s noise-only takes close the session: quiet pool, thruster,
bubbles or splashing, and pumps.\footnote{These takes were added to the S1 plan on 2026-11-09, before
any take was recorded; the speech and anchor takes and the bench were unchanged.} Projector-chain (E2)
and IR-grid takes are deferred to S2.

\begin{table}[t]
\caption{S1 take plan, in recording order (position \texttt{p3}, 3~m; thr.\ = thruster).}
\label{tab:takes}
\centering
\small
\setlength{\tabcolsep}{3.5pt}
\begin{tabular}{lll}
\hline
Take & Content & Length \\
\hline
\texttt{p3\_noise} & nothing played, thr.\ off & 1.0~min \\
\texttt{p3\_noise\_on} & nothing played, thr.\ on & 1.0~min \\
\texttt{p3\_sweeps} & 1 chirp, 4 ESS & 1.0~min \\
\texttt{p3\_speech} & 56 blocks, 112 chirps, 480 utt. & 66.5~min \\
\texttt{p3\_anchor} & 1 block, 2 chirps, 9 utt. & 1.2~min \\
\texttt{noise\_*} ($\times$4) & noise bank & 1.5~min each \\
\hline
\end{tabular}
\end{table}

\subsection{Safety}
\label{sec:safety}
The session requires written approval from the pool operator and an answer from the institution's
environmental health and safety (EH\&S) office on whether powered equipment and thrusters need its
sign-off; both are to be confirmed and logged for S1 \tbd. The battery-only rig keeps mains equipment
away from the pool. Nobody is in the water while the projector transmits. If anyone must
be, the received level measured at their location must stay at or below 150~dB re 1~$\mu$Pa (RMS over
speech-active frames), our own conservative limit. Nobody comes within 2~m of the running thruster,
which has a guard frame and an emergency stop. Because the projector can generate harmonics that the
drive LPF cannot remove, a raw 96~kHz capture is checked for energy in the 25--40~kHz band of the
vehicle's acoustic pinger; if any is found, the drive is lowered.

\subsection{Logging and metadata}
\label{sec:logging}
For every take, the session tool logs the measured range, projector and hydrophone depths, water
temperature, amplifier and recorder gains, pump state, pool activity, thruster unit and throttle, and
operator to \texttt{takes.jsonl} and to a printed run sheet. Water temperature is also logged at 0.5
and 1.5~m at the start and end of the session, with photographs of the rig; operators keep at least 2~m
from the array during takes, and positions are measured to a target of $\pm3$~cm. The import writes
one row per utterance and channel with the identity and alignment fields (session order, block,
channel, raw file and sample bounds, drift, chirp QC and flags); geometry, gains, temperatures and
acoustic covariates are merged from the session log afterwards. The S1 values of all logged quantities
are \tbd.

\subsection{Quality-control pass criteria}
\label{sec:qc}
A speech or anchor take passes when, on every recorder file, every chirp is found with a PSR of at
least 6~dB, no block is flagged, and no sample reaches full scale ($|x| \ge 0.999$). A take that fails
is recorded again before the rig is moved. Separately, each Daisy sidecar file is checked for dropped
frames; a single drop is not visible to the chirp QC, so a take with drops is redone only if many of its
blocks' drift estimates disagree with the take drift. The
per-take drift, minimum PSR and flagged blocks are reviewed after import. The S1 QC outcomes are \tbd.

\subsection{Per-utterance SNR and level stability}
\label{sec:snr}
Each utterance's SNR is measured from two recordings through the same chain and channel. Speech power
is the band power over speech-active samples (energy VAD) of the utterance's own thruster-off cut at the
same session, position and channel; ABBA guarantees that this cut exists and is adjacent in time. Noise
power comes from the position's noise-only take in the row's thruster state, over the whole take minus
1~s at each end, keeping only the 0.1~s frames whose 50--8000~Hz level lies within 6~dB below the
take's median frame level. SNR is reported over 500--8000~Hz, 50--8000~Hz and the octaves from 500 to
4000~Hz. The level drift of the session is the median, over the anchor's utterances, of the
500--8000~Hz speech-active level of the anchor cut minus that of the original cut; if its magnitude
exceeds 1.0~dB (a threshold fixed before any data), every cut of the session is flagged. This method
assumes that thruster noise is stationary between the noise-only take and the speech blocks. The flag is
reported and never used to filter a comparison. The measured S1 SNRs and drift are \tbd.

\subsection{References for intrusive metrics}
\label{sec:refs}
The real-data reference is $R_1 = s * h_{\mathrm{chain}}$, with $h_{\mathrm{chain}}$ measured once at
close range per source depth and only a broadband delay and gain fitted per utterance
\cite{yang2024realman}. S1 records no projector-chain measurement, so $R_1$ and
the metrics that depend on it are not available from S1 alone \tbd. The aligned $R_0$ cut is used for
the VAD of level normalization.

\subsection{Primary endpoint and secondary analyses}
\label{sec:endpoint}
The S1 scope and endpoint were signed off in the project repository on 2026-11-09, before any pool
recording existed; the external pre-registration (OSF) is to be filed before v0-test is scored \tbd.
This amendment replaces, for S1 only, the originally planned primary endpoint, the paired
$\Delta\mathrm{WER}$ of C2 minus C0 pooled over all v0 positions and thruster states. On S1, C2 minus
C0 is secondary; for the pooled v0 after session S2, the plan's endpoints are unchanged.

\textit{Primary endpoint.} The paired $\Delta\mathrm{WER}$ of C8 (the frozen adapter of
Section~\ref{sec:c8}) minus C0 on the S1 v0-test cuts of the speech take on the scored channel
\texttt{ch0} (200 utterances in each of the two thruster states, 400 cuts), with the frozen guarded
decoding (Section~\ref{sec:decode}) and statistics (Section~\ref{sec:stats}). The v0-dev cuts and the
drift-anchor repeats are excluded; the anchor serves only the level-drift check. The endpoint is reported overall and by thruster state. Result: \tbd.

\textit{Secondary analyses.} (i) Each SE front end (C1, C2, C10) minus C0, and each SE front end
minus C8 (paired $\Delta\mathrm{WER}$, no directional prediction), scored once a front-end inference
script for real cuts exists. (ii) Flag rates per system and WER without flagged
utterances on the paired subset. (iii) WER by speaker sex against the two-band SNR. (iv) The anchor
level-drift flag. All other analyses are exploratory and labelled as such.

\textit{Scope.} With a single position, range effects cannot be tested on S1, and any result is
scoped to one pool, one position and one session. S1 has no test--retest repeat, so the noise floor
below which no difference is interpreted is not available from S1 \tbd.

\subsection{Post-session channel characterisation}
\label{sec:post}
On the day of the session, a CPU-only analysis is planned for the non-speech takes. Levels are
uncalibrated and are compared only within a channel.

\textit{Sweep IRs.} The clock drift and offset of each sweep recording are estimated from the sweeps
themselves on the recorder's sync channel, and one IR per sweep and hydrophone is obtained by
regularized deconvolution against the regenerated excitation. From the IRs we report sweep coherence
and between-sweep variance; the energy decay curve (EDC) by Schroeder integration with Lundeby
truncation \cite{schroeder1965,lundeby1995}; broadband and octave-band
(250~Hz--8~kHz) early decay time (EDT, 0 to $-10$~dB), T20 ($-5$ to $-25$~dB) and T30 ($-5$ to
$-35$~dB); a single- or two-slope decay fit; the clarity index C50 (energy before and after 50~ms);
the direct-to-reverberant ratio (DRR); and the IR SNR per band. DRR uses a direct-path gate of 0.7~ms
around the IR peak and is reported as missing where the gate rule fails, that is, where the gate reaches
the first reflection or holds fewer than two periods of the lowest analysed frequency. Arrivals more
than 20~dB below the strongest are not counted as reflections. S1 values: \tbd.

\textit{Noise spectra.} For each noise take, with the same 1~s trim and 6~dB median gate as the SNR
method, we report the mean power spectral density (Hann frames of about 85~ms, 50\% overlap) and
one-third-octave band levels in dBFS, with their frame-to-frame spread as a measure of stationarity,
and each condition relative to the quiet-pool take. S1 values: \tbd.

% ---------------------------------------------------------------------------------------------
\section{Results}
% TODO: S1 primary endpoint (C8 - C0, overall and by thruster state) and secondary analyses; simulation
% results (labelled as simulation) as in the abstract.
\textit{TODO: Results (real data \tbd; simulation results labelled as such).}

\section{Discussion}
% TODO
\textit{TODO: Discussion.}

\section{Conclusion}
% TODO
\textit{TODO: Conclusion.}

% Bibliography copied from paper/oceans2027_abstract.tex; entries not yet cited will be cited by the
% Introduction and Related Work. Do not add entries without verifying them.
\begin{thebibliography}{99}
\bibitem{nissen2017} I.~Nissen and E.~Kuhnt-Matth\'e, ``Verbesserung der Sprachverst\"andlichkeit bei
analoger Unterwasserkommunikation durch Nutzung von digitalen Anfangs- und Endmarken,'' in
\emph{Proc. DAGA}, Kiel, 2017.
\bibitem{roehrdanz2025} F.~R\"ohrdanz \emph{et al.}, ``Data generation for training of neural networks
dealing with underwater speech transmission,'' in \emph{Proc. DAS|DAGA}, Copenhagen, 2025.
\bibitem{roehrdanz2026} F.~R\"ohrdanz \emph{et al.}, ``Neural-network based speech dereverberation for
underwater applications,'' in \emph{Proc. DAGA}, 2026, no.~144.
\bibitem{allen1979} J. B. Allen and D. A. Berkley, ``Image method for efficiently simulating small-room acoustics,'' \emph{J. Acoust. Soc. Am.}, vol. 65, no. 4, pp. 943--950, Apr. 1979, doi: 10.1121/1.382599.
\bibitem{brekhovskikh1980} L. M. Brekhovskikh, \emph{Waves in Layered Media}, 2nd ed., R. T. Beyer, Trans. New York, NY, USA: Academic Press, 1980, ser. Applied Mathematics and Mechanics, vol. 16.
\bibitem{minnaert1933} M. Minnaert, ``On musical air-bubbles and the sounds of running water,'' \emph{London, Edinburgh, Dublin Philos. Mag. J. Sci.}, vol. 16, no. 104, pp. 235--248, Aug. 1933, doi: 10.1080/14786443309462277.
\bibitem{panayotov2015} V.~Panayotov, G.~Chen, D.~Povey, and S.~Khudanpur, ``LibriSpeech: an ASR corpus
based on public domain audio books,'' in \emph{Proc. ICASSP}, 2015, pp.~5206--5210.
\bibitem{nakatani2010} T.~Nakatani, T.~Yoshioka, K.~Kinoshita, M.~Miyoshi, and B.-H.~Juang, ``Speech dereverberation based on variance-normalized delayed linear prediction,'' \emph{IEEE Trans. Audio, Speech, Lang. Process.}, vol.~18, no.~7, pp.~1717--1731, Sep. 2010, doi: 10.1109/TASL.2010.2052251.
\bibitem{drude2018} L.~Drude, J.~Heymann, C.~Boeddeker, and R.~Haeb-Umbach, ``NARA-WPE: A Python package for weighted prediction error dereverberation in Numpy and Tensorflow for online and offline processing,'' in \emph{Proc. 13th ITG Fachtagung Sprachkommunikation (ITG-Fachbericht 282)}, Oldenburg, Germany, Oct. 2018, pp.~1--5.
\bibitem{fu2021} S.-W.~Fu, C.~Yu, T.-A.~Hsieh, P.~Plantinga, M.~Ravanelli, X.~Lu, and Y.~Tsao, ``MetricGAN+: An improved version of MetricGAN for speech enhancement,'' in \emph{Proc. Interspeech}, 2021, pp.~201--205, doi: 10.21437/Interspeech.2021-599.
\bibitem{subakan2021} C.~Subakan, M.~Ravanelli, S.~Cornell, M.~Bronzi, and J.~Zhong, ``Attention is all you need in speech separation,'' in \emph{Proc. IEEE Int. Conf. Acoust., Speech, Signal Process. (ICASSP)}, Toronto, ON, Canada, 2021, pp.~21--25, doi: 10.1109/ICASSP39728.2021.9413901.
\bibitem{wichern2019} G.~Wichern, J.~Antognini, M.~Flynn, L.~R.~Zhu, E.~McQuinn, D.~Crow, E.~Manilow, and J.~Le~Roux, ``WHAM!: Extending speech separation to noisy environments,'' in \emph{Proc. Interspeech}, 2019, pp.~1368--1372, doi: 10.21437/Interspeech.2019-2821.
\bibitem{maciejewski2020} M.~Maciejewski, G.~Wichern, E.~McQuinn, and J.~Le~Roux, ``WHAMR!: Noisy and reverberant single-channel speech separation,'' in \emph{Proc. IEEE Int. Conf. Acoust., Speech, Signal Process. (ICASSP)}, Barcelona, Spain, 2020, pp.~696--700, doi: 10.1109/ICASSP40776.2020.9053327.
\bibitem{schroeter2023} H.~Schr\"oter, A.~N. Escalante-B., T.~Rosenkranz, and A.~Maier,
``DeepFilterNet: Perceptually motivated real-time speech enhancement,'' in \emph{Proc. Interspeech},
2023, pp.~2008--2009.
\bibitem{radford2023} A.~Radford \emph{et al.}, ``Robust speech recognition via large-scale weak supervision,'' in \emph{Proc. ICML}, 2023,
pp.~28492--28518.
\bibitem{hu2022} E.~J. Hu \emph{et al.}, ``LoRA: Low-rank adaptation of large language models,'' in
\emph{Proc. ICLR}, 2022.
\bibitem{rong2025} Y.~Rong, P.~Chen, and Z.~Nan, ``Underwater diver communication system based on speech/text conversion,'' in \emph{Proc. OCEANS 2025 Brest}, Brest, France, Jun. 2025, pp.~1--5, doi: 10.1109/OCEANS58557.2025.11104740.
\bibitem{jensen2016} J.~Jensen and C.~H. Taal, ``An algorithm for predicting the intelligibility of
speech masked by modulated noise maskers,'' \emph{IEEE/ACM Trans. Audio, Speech, Lang. Process.},
vol.~24, no.~11, pp.~2009--2022, 2016.
\bibitem{farina2000} A.~Farina, ``Simultaneous measurement of impulse response and distortion with a swept-sine technique,'' presented at the \emph{108th AES Convention}, Paris, France, Feb. 19--22, 2000, Preprint 5093.
\bibitem{yang2024realman} B.~Yang, C.~Quan, Y.~Wang, P.~Wang, Y.~Yang, Y.~Fang, N.~Shao, H.~Bu, X.~Xu, and X.~Li, ``RealMAN: A real-recorded and annotated microphone array dataset for dynamic speech enhancement and localization,'' in \emph{Advances in Neural Information Processing Systems}, vol.~37, 2024, pp.~105997--106019, doi: 10.52202/079017-3363.
\bibitem{schroeder1965} M.~R. Schroeder, ``New method of measuring reverberation time,'' \emph{J. Acoust. Soc. Am.}, vol.~37, no.~3, pp.~409--412, 1965, doi: 10.1121/1.1909343.
\bibitem{lundeby1995} A.~Lundeby, T.~E. Vigran, H.~Bietz, and M.~Vorl\"ander, ``Uncertainties of measurements in room acoustics,'' \emph{Acustica}, vol.~81, no.~4, pp.~344--355, 1995.
\bibitem{iwamoto2022} K.~Iwamoto \emph{et al.}, ``How bad are artifacts?: Analyzing the impact of speech enhancement errors on ASR,'' in
\emph{Proc. Interspeech}, 2022, pp.~5418--5422.
\end{thebibliography}

\end{document}
